\section*{A brief History of Nuclear Power, from its origins to present day}
The 1930s and 1940s witnessed unprecedented acceleration in nuclear science. Following James Chadwick's 1932 discovery of the neutron \cite{PLACEHOLDER}, Otto Hahn and Fritz Strassmann demonstrated nuclear fission in 1938 by producing barium through uranium neutron bombardment \cite{PLACEHOLDER}. This breakthrough enabled Enrico Fermi's team to construct Chicago Pile-1 in 1942 - the first artificial nuclear reactor capable of sustaining and controlling a chain reaction \cite{PLACEHOLDER} - just four years after the fission discovery.

While initial advancements were primarily motivated by military applications during World War II, culminating in the atomic bombings of Hiroshima (August 6, 1945) and Nagasaki (August 9, 1945), the scientific community quickly recognized nuclear technology's peaceful potential. President Eisenhower's 1953 "Atoms for Peace" address to the United Nations \cite{PLACEHOLDER} established the foundation for international nuclear cooperation, leading to the 1957 creation of both the International Atomic Energy Agency (IAEA) and the European Atomic Energy Community (EURATOM). These institutions were tasked with promoting peaceful nuclear applications while serving as independent authorities for safety regulation and non-proliferation oversight.

The 1950s and 1960s marked a period of rapid nuclear reactor development, driven by dual civilian and military imperatives. This era saw several key milestones in nuclear power generation: the Experimental Breeder Reactor I (EBR-I) in Idaho (USA) first produced electricity in 1951 (powering four light bulbs) \cite{PLACEHOLDER}, followed by the Obninsk Nuclear Power Plant (USSR) in 1954, which became the first to feed approximately 5 MWe into a local grid \cite{PLACEHOLDER}. The world's first commercial-scale nuclear power station, Calder Hall-1 (UK), began operation in 1956, connecting to the national grid with an initial capacity of 50 MWe \cite{PLACEHOLDER}.

Parallel technological diversification occurred during this period. The United Kingdom developed gas-cooled Magnox reactors, optimized for both electricity production and plutonium generation. Concurrently, the Soviet Union initiated development of the graphite-moderated, water-cooled RBMK design through experimental work at Obninsk. Canada's nuclear program focused on pressurized heavy water reactor (PHWR) technology, culminating in the 1962 Rolphton demonstration reactor and the first commercial CANDU reactor at Douglas Point in 1968 \cite{PLACEHOLDER}.

In reactor cooling technologies, the United States pioneered multiple approaches: the EBR-I employed liquid metal (NaK) cooling, while the naval propulsion program led to light water reactor (LWR) development with the 1954 launch of USS Nautilus. By 1957, both Boiling Water Reactor (BWR) and Pressurized Water Reactor (PWR) designs were adapted for civilian applications, with the Shippingport Reactor (Pennsylvania) becoming the first full-scale civilian PWR connected to the grid. The Soviet VVER ("Water-Water Energetic Reactor") program advanced simultaneously, deploying the 210 MWe VVER-210 at Novovoronezh Nuclear Power Plant in 1964 \cite{PLACEHOLDER}.

The 1960s and 1970s witnessed significant global expansion of nuclear power generation. During this period, numerous countries commissioned their first nuclear reactors: Argentina (Atucha-1, 1974), Armenia (Armenian-1, 1976), Belgium (BR-3, 1962), Bulgaria (1974), Italy (Latina, 1963), Japan (1963), Sweden (1964), Switzerland (1969), and Slovakia (1972) \cite{PLACEHOLDER}.

The 1973 oil crisis accelerated nuclear adoption as an energy security strategy. France implemented the Messmer Plan, an ambitious program that resulted in the construction of 43 reactors between 1977 and 1986 \cite{PLACEHOLDER}. Similarly, Japan significantly expanded its nuclear capacity to reduce dependence on imported oil, establishing nuclear energy as a cornerstone of its national energy strategy \cite{PLACEHOLDER}.

The nuclear expansion continued through the 1980s until the 1986 Chernobyl disaster, which constituted a critical inflection point for global nuclear development. The accident prompted widespread safety reviews and regulatory updates, significantly slowing deployment rates \cite{PLACEHOLDER}. Public opposition intensified particularly in Western Europe, leading several countries to halt or reduce their nuclear programs. Italy's 1987 referendum resulted in the complete phase-out of nuclear power, while other European nations experienced similar setbacks \cite{PLACEHOLDER}.

By the 1990s, multiple countries had suspended nuclear deployment entirely, including Germany and Switzerland. France, previously the most active builder, connected only 10 reactors between 1990-1999, while the United States and United Kingdom deployed 4 and 1 reactors respectively during the same period \cite{PLACEHOLDER}. Italy completed its nuclear phase-out by 1990 with the permanent shutdown of all operating reactors. Following the post-Soviet economic transition, Russia resumed its nuclear program and has since maintained continuous expansion of its reactor fleet \cite{PLACEHOLDER}.

A nuclear renaissance emerged in the early 2000s, driven by concerns over climate change, energy security, and rising fossil fuel prices. This period saw the development of advanced reactor designs including the European Pressurized Reactor (EPR) and AP1000 \cite{PLACEHOLDER}. However, the 2011 Fukushima Daiichi accident constituted another significant setback for nuclear expansion in Western nations. Projects faced substantial delays, cost overruns, and public opposition, with France's Flamanville-3 EPR serving as a particularly illustrative case: experiencing budget increases to eight times original estimates and construction timelines tripled from initial projections \cite{PLACEHOLDER}. Germany accelerated its nuclear phase-out policy, while Italy reaffirmed its anti-nuclear position through a popular referendum \cite{PLACEHOLDER}.

Contrastingly, Eastern nations pursued markedly different nuclear trajectories. China, having initiated its nuclear program in the 1990s with Qinshan-1 (connected in 1991), has since maintained aggressive expansion. Over the past three decades, China has brought 56 additional reactors online and currently has 29 under construction, positioning itself to surpass both France and the United States in total reactor capacity within coming decades \cite{PLACEHOLDER}. This expansion reflects a strategic commitment to nuclear energy as part of China's broader energy and climate objectives.

The COVID-19 pandemic in 2020 exposed critical vulnerabilities in global economic systems, while Russia's 2022 invasion of Ukraine further revealed Europe's energy security dependencies \cite{PLACEHOLDER}. Concurrently, decarbonization imperatives to combat climate change have prompted Western nations to reconsider nuclear energy's dual role: providing low-carbon, reliable baseload power while enhancing energy independence and system resilience \cite{EU2050ClimateStrategy}.

This strategic reassessment is evident in recent policy initiatives. The European Union's classification of nuclear power as a sustainable investment under its taxonomy for climate mitigation \cite{EUTaxonomy}, alongside programs like the U.S. Department of Energy's Advanced Reactor Demonstration Program \cite{PLACEHOLDER}, demonstrates renewed interest in nuclear technologies. The European energy sector is consequently undergoing a significant transformation, balancing decarbonization targets with energy security requirements through diversified generation portfolios that include advanced nuclear solutions \cite{EU2050ClimateStrategy}.

While the political decision to include nuclear power is fundamentally a strategic choice, selecting specific reactor designs involves complex trade-offs across technical, economic, regulatory, social, environmental, and geopolitical dimensions \cite{PLACEHOLDER}. National energy needs vary significantly: some countries require large baseload capacity to reduce import dependence, while others prioritize smaller, more flexible solutions to integrate with existing renewable systems \cite{PLACEHOLDER}.

The current nuclear landscape presents particular challenges for decision-makers, offering both established Generation III+ designs and emerging Small Modular Reactors (SMRs) as potential options. Given that nuclear commitments represent decade-long investments with 60-year operational lifespans \cite{PLACEHOLDER} and significant financial implications, and considering the multiple influencing factors like grid infrastructure, public acceptance, and supply chain considerations, policymakers require structured decision-support tools to evaluate these complex trade-offs. This study therefore employs Multi-Criteria Decision Analysis (MCDA), specifically the Multi-Attribute Value Theory (MAVT) approach, presented in the following paragraphs.

\subsection*{MCDA and MAVT as a Decision Support Tool}
Multi-Criteria Decision Analysis (MCDA) provides a structured, transparent framework for addressing complex decision problems where trade-offs between competing objectives are inevitable. Common to any decision making process is a systematic process that involves: defining the decision problem, identifying relevant objectives, establishing alternative decision options, and quantifying the performance of each option against these objectives. \\

Among these methods, Multi-Attribute Value Theory (MAVT) employs a hierarchical objective structure and value functions to convert both qualitative and quantitative data into comparable units. By assigning importance weights to each attribute, MAVT aggregates these measurements into a single composite score for each option. This approach is particularly effective for ranking alternatives and providing decision-makers with a comprehensive comparison of available options. \\

\textcolor{orange}{how much more do I have to explain here? I guess that I will explain the way we apply weights when we do it, the way swing elicitation works when we apply it etc etc...}